\documentclass[11pt]{article}

% --- Codificación e idioma ---
\usepackage[T1]{fontenc}
\usepackage[utf8]{inputenc}
\usepackage[spanish,english]{babel}

% --- Matemáticas ---
\usepackage{amsmath,amsthm,amssymb}

% --- Tablas y formato ---
\usepackage{float}
\usepackage{booktabs}
\usepackage[normalem]{ulem}
\usepackage[a4paper,margin=2.5cm]{geometry}
\usepackage[hidelinks]{hyperref}

% --- Numeración ---
\numberwithin{equation}{section}
\numberwithin{figure}{section}
\numberwithin{table}{section}

% --- Entornos ---
\theoremstyle{definition}
\newtheorem*{definition}{Definición}
\newtheorem*{restriccion}{Restricción}

% --- Nombres en español ---
\addto\shorthandsspanish{\spanishdeactivate{~<>}}

\title{Algoritmo de asignación de horarios estudiantiles\\
\large Entrega 1: Definición formal y estrategia de datos}
\author{Cesar Santiago Canro Cabra \and Samuel Enrique Giraldo Sabogal}
\date{\today}

\begin{document}
\selectlanguage{spanish}

\maketitle

\begin{abstract}
En este documento se presenta el análisis y el diseño formal del problema de
asignación de horarios estudiantiles en la Pontificia Universidad Javeriana,
así como la estrategia de datos propuesta para probar y evaluar el algoritmo
de solución. Se formalizan las entradas, las precondiciones, las restricciones
de validez, las preferencias de los estudiantes y la salida esperada. En esta
etapa no se define ningún algoritmo de solución; esa decisión se desarrollará
en la siguiente entrega.
\end{abstract}

\tableofcontents
\newpage

\part{Definición formal del problema}

\section{Análisis del problema}

Se dispone de un conjunto de estudiantes que deben construir su horario
académico a partir de las asignaturas que desean cursar. Cada asignatura puede
ofrecer uno o varios grupos, y cada grupo se encuentra caracterizado por un
profesor, una o varias franjas horarias, un salón y una capacidad máxima de
estudiantes.

El problema consiste en determinar, para cada estudiante, qué grupos deben ser
asignados de manera que el horario resultante sea válido. Una asignación se
considera válida cuando cumple simultáneamente las siguientes condiciones:

\begin{enumerate}
  \item un estudiante no puede estar asignado a dos grupos cuyas franjas
        horarias se intersecten;
  \item un estudiante puede pertenecer, como máximo, a un grupo por cada
        asignatura solicitada;
  \item la cantidad de estudiantes asignados a un grupo no puede superar la
        capacidad máxima de dicho grupo;
  \item un estudiante solamente puede ser asignado a grupos correspondientes
        a asignaturas que haya solicitado.
\end{enumerate}

Satisfacer estas restricciones determina la \emph{validez} de una asignación,
pero no necesariamente una única solución. Para un mismo estudiante pueden
existir varios horarios válidos. Por esta razón, además de garantizar la
validez, se desea contar con información que permita evaluar la calidad del
horario obtenido.

Los principales criterios de calidad considerados son la prioridad de las
asignaturas, la preferencia por determinados profesores, la organización del
tiempo libre, la cercanía entre salones de clases consecutivas y la posibilidad
de coincidir con determinados compañeros.

Por tanto, en el problema se distinguen dos tipos de condiciones:

\begin{itemize}
  \item \textbf{Restricciones de validez:} condiciones obligatorias que todo
        horario generado debe satisfacer.
  \item \textbf{Criterios de calidad:} características que permiten evaluar y
        comparar horarios que ya son válidos.
\end{itemize}

El problema puede resumirse como la búsqueda de una asignación de grupos para
cada estudiante que produzca horarios válidos y cuya calidad pueda ser
evaluada mediante las prioridades y preferencias recibidas como entrada. En
esta entrega no se establece aún el procedimiento algorítmico para encontrar
la asignación; el análisis determina únicamente qué representa una solución
correcta y qué información resulta relevante para evaluarla.

\section{Diseño del problema}

El diseño establece el contrato formal del problema con independencia del
algoritmo que posteriormente se utilice. Dicho contrato se organiza en cuatro
componentes:

\begin{enumerate}
  \item \textbf{Entradas:} estudiantes, asignaturas, grupos, profesores,
        salones, franjas horarias, solicitudes, prioridades y preferencias.
  \item \textbf{Precondiciones:} propiedades que deben satisfacer los datos
        recibidos antes de realizar la asignación.
  \item \textbf{Salida:} función que asocia a cada estudiante el conjunto de
        grupos que conforma su horario.
  \item \textbf{Postcondiciones:} restricciones que debe cumplir la salida
        para ser considerada válida.
\end{enumerate}

Una vez verificada la validez de la salida, se calculan métricas de calidad que
permiten comparar distintas asignaciones válidas. Las secciones siguientes
desarrollan cada componente del contrato.

\section{Especificación formal de las entradas}

\subsection{Conjuntos}

Se definen los siguientes conjuntos finitos:

\begin{itemize}
  \item $E = \{e_1,e_2,\ldots,e_n\}$: conjunto de estudiantes.
  \item $A = \{a_1,a_2,\ldots,a_m\}$: conjunto de asignaturas disponibles.
  \item $G = \{g_1,g_2,\ldots,g_k\}$: conjunto de grupos o secciones.
  \item $P = \{p_1,p_2,\ldots,p_r\}$: conjunto de profesores.
  \item $S = \{s_1,s_2,\ldots,s_q\}$: conjunto de salones.
  \item $D = \{1,2,3,4,5\}$: conjunto ordenado de días académicos de la
        semana.
  \item $T$: conjunto ordenado de instantes de tiempo considerados durante
        un día académico.
  \item $F = \{f_1,f_2,\ldots,f_t\}$: conjunto de franjas horarias
        disponibles durante la semana.
\end{itemize}

Cada franja $f\in F$ se representa mediante una terna:

\[
f=\bigl(d(f),t_{\mathrm{inicio}}(f),t_{\mathrm{fin}}(f)\bigr)
\in D\times T\times T,
\]

con $t_{\mathrm{inicio}}(f)<t_{\mathrm{fin}}(f)$. Esta representación permite
ordenar temporalmente las clases, calcular los tiempos libres y determinar
cuándo dos franjas se superponen.

Se define la relación de solapamiento $\bowtie$ sobre $F$. Para
$f_i,f_j\in F$:

\[
f_i\bowtie f_j
\iff
d(f_i)=d(f_j)
\land
t_{\mathrm{inicio}}(f_i)<t_{\mathrm{fin}}(f_j)
\land
t_{\mathrm{inicio}}(f_j)<t_{\mathrm{fin}}(f_i).
\]

Dos clases son consecutivas cuando ocurren el mismo día, no existe otra clase
asignada entre ellas y la hora de inicio de la segunda no es anterior a la
hora de finalización de la primera. Esta relación se utilizará posteriormente
para calcular tiempos libres y distancias entre salones consecutivos.

\subsection{Funciones sobre los grupos}

Cada grupo $g\in G$ pertenece exactamente a una asignatura. Se definen las
siguientes funciones:

\[
\begin{aligned}
  \operatorname{asignatura} &: G \rightarrow A,\\
  \operatorname{profesor}   &: G \rightarrow P,\\
  \operatorname{horario}    &: G \rightarrow 2^F,\\
  \operatorname{salon}      &: G \rightarrow S,\\
  \operatorname{capacidad}  &: G \rightarrow \mathbb{N}.
\end{aligned}
\]

Para cada grupo $g$, $\operatorname{asignatura}(g)$ indica la asignatura a la
que pertenece; $\operatorname{profesor}(g)$, el profesor encargado;
$\operatorname{horario}(g)$, el conjunto de franjas que ocupa;
$\operatorname{salon}(g)$, el salón donde se desarrolla; y
$\operatorname{capacidad}(g)$, el número máximo de estudiantes que puede
recibir.

\subsection{Información de los estudiantes}

Para cada estudiante $e\in E$ se define el conjunto de asignaturas solicitadas:

\[
\operatorname{Solicitudes}(e)\subseteq A.
\]

También se define la función de prioridad:

\[
\operatorname{prioridad}:E\times A\rightarrow\{1,2,\ldots,K\},
\]

donde $\operatorname{prioridad}(e,a)$ representa la importancia que tiene para
el estudiante $e$ cursar la asignatura $a$. Un valor mayor indica una prioridad
mayor.

\subsection{Información de los grupos}

Para cada grupo $g\in G$ se conoce:

\[
\operatorname{asignatura}(g)\in A,\quad
\operatorname{profesor}(g)\in P,\quad
\operatorname{horario}(g)\subseteq F,
\]

\[
\operatorname{salon}(g)\in S,\qquad
\operatorname{capacidad}(g)\in\mathbb{N}.
\]

\section{Preferencias de los estudiantes}

\subsection{Preferencia por profesor}

Se define la función:

\[
\operatorname{preferenciaProfesor}:E\times P
\rightarrow\{0,1,2,3,4,5\},
\]

donde $\operatorname{preferenciaProfesor}(e,p)$ representa el nivel de
preferencia del estudiante $e$ por recibir clase con el profesor $p$, según la
siguiente escala:

\begin{table}[H]
\centering
\caption{Escala de preferencia por profesor.}
\begin{tabular}{cl}
\toprule
\textbf{Valor} & \textbf{Significado} \\
\midrule
0 & Sin preferencia o profesor desconocido \\
1 & Preferencia muy baja \\
2 & Preferencia baja \\
3 & Preferencia media \\
4 & Preferencia alta \\
5 & Preferencia muy alta \\
\bottomrule
\end{tabular}
\end{table}

\subsection{Preferencias de horario}

Para cada estudiante $e\in E$ se define:

\[
\operatorname{PreferenciasHorario}(e)=
\left(h_{\min}(e),h_{\max}(e),c(e),l(e)\right),
\]

donde:

\begin{itemize}
  \item $h_{\min}(e)\in T$ es la hora mínima deseada para iniciar clases;
  \item $h_{\max}(e)\in T$ es la hora máxima deseada para terminar clases;
  \item $c(e)\in\{0,1\}$ expresa la preferencia por un horario compacto
        ($1$ si prefiere evitar huecos y $0$ en caso contrario);
  \item $l(e)\subseteq D$ es el conjunto de días de la semana en
        los que el estudiante prefiere no tener clases.
\end{itemize}

\subsection{Cercanía entre salones}

Se define la función:

\[
\operatorname{distancia}:S\times S\rightarrow\mathbb{R}_{\geq 0},
\]

donde $\operatorname{distancia}(s_i,s_j)$ representa la distancia física en
metros entre los salones $s_i$ y $s_j$. Para una instancia real se pueden
usar mediciones del campus. La estrategia de datos describe cómo simular
distancias para los experimentos.

\subsection{Preferencias sociales y académicas}

Para cada estudiante $e\in E$ se define:

\[
\operatorname{Compañeros}(e)\subseteq E\setminus\{e\},
\]

que representa el conjunto de estudiantes con los cuales $e$ preferiría
coincidir en uno o más grupos.

\section{Precondiciones}

Las entradas deben satisfacer las siguientes propiedades antes de ejecutar
cualquier algoritmo de asignación:

\begin{enumerate}
  \item Los conjuntos $E$, $A$, $G$, $P$, $S$, $D$, $T$ y $F$ son finitos.
  \item Para toda franja $f\in F$:
        \[
        d(f)\in D,\qquad
        t_{\mathrm{inicio}}(f),t_{\mathrm{fin}}(f)\in T,
        \]
        \[
        t_{\mathrm{inicio}}(f)<t_{\mathrm{fin}}(f).
        \]
  \item Para todo $g\in G$:
        \[
        \operatorname{asignatura}(g)\in A,\quad
        \operatorname{profesor}(g)\in P,\quad
        \operatorname{horario}(g)\subseteq F,
        \]
        \[
        \operatorname{salon}(g)\in S,\qquad
        \operatorname{capacidad}(g)>0.
        \]
  \item Para todo $e\in E$:
        \[
        \emptyset\neq\operatorname{Solicitudes}(e)\subseteq A.
        \]
  \item Para todo par $s_i,s_j\in S$:
        \[
        \operatorname{distancia}(s_i,s_j)\geq 0,
        \]
        \[
        \operatorname{distancia}(s_i,s_i)=0,
        \]
        \[
        \operatorname{distancia}(s_i,s_j)
        =\operatorname{distancia}(s_j,s_i).
        \]
  \item Para todo $e\in E$:
        \[
        \operatorname{Compañeros}(e)\subseteq E\setminus\{e\}.
        \]
\end{enumerate}

\section{Salida formal y postcondiciones}

La salida principal del problema es la función:

\[
H:E\rightarrow 2^G,
\]

donde $H(e)$ representa el conjunto de grupos asignados al estudiante $e$.
La salida completa puede expresarse como:

\[
H=\left\{\bigl(e,H(e)\bigr)\mid e\in E\right\}.
\]

Cada conjunto $H(e)$ debe cumplir las siguientes postcondiciones.

\begin{restriccion}[Asignación de asignaturas solicitadas]
Para todo estudiante $e\in E$ y todo grupo $g\in H(e)$:
\[
\operatorname{asignatura}(g)\in\operatorname{Solicitudes}(e).
\]
\end{restriccion}

\begin{restriccion}[Ausencia de cruces de horario]
Para todo estudiante $e\in E$, todo par de grupos distintos
$g_i,g_j\in H(e)$ y toda pareja de franjas
$f_i\in\operatorname{horario}(g_i)$,
$f_j\in\operatorname{horario}(g_j)$:
\[
\neg\bigl(f_i\bowtie f_j\bigr).
\]
\end{restriccion}

\begin{restriccion}[Una sola sección por asignatura]
Para todo estudiante $e\in E$ y todo par de grupos distintos
$g_i,g_j\in H(e)$:
\[
\operatorname{asignatura}(g_i)\neq\operatorname{asignatura}(g_j).
\]
\end{restriccion}

\begin{restriccion}[Respeto de la capacidad]
Para todo grupo $g\in G$:
\[
\left|\{e\in E\mid g\in H(e)\}\right|
\leq\operatorname{capacidad}(g).
\]
\end{restriccion}

\section{Métricas de calidad}

Una vez comprobado que $H$ satisface las postcondiciones, se asocia a cada
estudiante un vector de métricas:

\[
M_e=\left(m_1(e),m_2(e),\ldots,m_z(e)\right).
\]

La salida extendida se define como:

\[
\operatorname{Resultado}(e)=\bigl(H(e),M_e\bigr).
\]

Las métricas pueden incluir:

\begin{itemize}
  \item número de asignaturas solicitadas y asignadas;
  \item satisfacción de las asignaturas de mayor prioridad;
  \item cantidad total de horas libres entre clases;
  \item distancia total recorrida entre clases consecutivas;
  \item satisfacción promedio con los profesores asignados;
  \item coincidencias con compañeros preferidos.
\end{itemize}

Para formalizar la cobertura, se define primero el conjunto de asignaturas
asignadas al estudiante $e$:

\[
\operatorname{Asignadas}(e)=
\left\{a\in\operatorname{Solicitudes}(e)\,\middle|\,
\exists g\in H(e):\operatorname{asignatura}(g)=a\right\}.
\]

La cobertura individual es:

\[
\operatorname{Cobertura}(e)=
\frac{|\operatorname{Asignadas}(e)|}
     {|\operatorname{Solicitudes}(e)|}.
\]

La cobertura global es:

\[
\operatorname{CoberturaGlobal}=
\frac{\displaystyle\sum_{e\in E}|\operatorname{Asignadas}(e)|}
     {\displaystyle\sum_{e\in E}|\operatorname{Solicitudes}(e)|}.
\]

Estas métricas permiten comparar asignaciones válidas sin afirmar todavía que
una de ellas es óptima respecto de una función objetivo no definida.

\subsection{Ejemplo conceptual de salida}

\begin{table}[H]
\centering
\caption{Ejemplo de horario asignado al estudiante $e_1$.}
\begin{tabular}{ll}
\toprule
\textbf{Asignatura} & \textbf{Grupo asignado} \\
\midrule
Análisis de Algoritmos & G03 \\
Bases de Datos         & G07 \\
Computación Móvil      & G11 \\
Optimización           & G15 \\
\bottomrule
\end{tabular}
\end{table}

\begin{table}[H]
\centering
\caption{Métricas del horario asignado al estudiante $e_1$.}
\begin{tabular}{lr}
\toprule
\textbf{Métrica} & \textbf{Valor} \\
\midrule
Materias solicitadas        & 4 \\
Materias asignadas          & 4 \\
Cruces de horario           & 0 \\
Tiempo libre total          & 6 horas \\
Distancia total recorrida   & 850 metros \\
Satisfacción con profesores & 4.25 / 5 \\
Prioridades satisfechas     & 100\% \\
\bottomrule
\end{tabular}
\end{table}

\part{Estrategia de datos}

\section{Formato de almacenamiento CSV}
\label{sec:csv}

Cada instancia se almacena en una carpeta con once archivos CSV. Una fila
representa una entidad o una relación. Por ejemplo, si un grupo se reúne
dos días, sus datos fijos se escriben una sola vez en \texttt{grupos.csv}
y las dos sesiones se escriben como filas en \texttt{horarios.csv}.
Así se pueden añadir sesiones o solicitudes sin modificar el encabezado
ni repetir el cupo, el profesor o el salón de un grupo.

\subsection{Convenciones y diccionario de datos}

Todos los CSV tienen encabezado obligatorio, columnas en el orden documentado,
separador coma y codificación UTF-8. Las columnas de un registro deben estar
completas. Los identificadores se tratan como texto y deben conservarse
al importar o editar los archivos. Los campos numéricos son enteros decimales
sin relleno con ceros. Las horas se expresan en minutos desde medianoche
(\texttt{480} equivale a 08:00). Los días van de 1 (lunes) a 5 (viernes).
Las métricas y los grupos asignados se guardan como salida JSON en una ruta
separada de las entradas.

\begin{table}[H]
\centering
\caption{Catálogos y entidades principales del esquema CSV.}
\small
\begin{tabular}{@{}p{0.27\linewidth}p{0.69\linewidth}@{}}
\toprule
\textbf{Archivo} & \textbf{Columnas en orden y clave primaria} \\
\midrule
\texttt{asignaturas.csv} &
\texttt{asignatura\_id}. Clave: asignatura. \\
\texttt{profesores.csv} &
\texttt{profesor\_id}. Clave: profesor. \\
\texttt{salones.csv} &
\texttt{salon\_id}. Clave: salón. \\
\texttt{estudiantes.csv} &
\texttt{estudiante\_id, inicio\_min, fin\_max, compacto}.\newline
Clave: estudiante. \texttt{compacto} vale 0 o 1. \\
\texttt{grupos.csv} &
\texttt{grupo\_id, asignatura\_id, profesor\_id,}\newline
\texttt{salon\_id, capacidad}. Clave: grupo. Asignatura,
profesor y salón deben existir en sus catálogos. \\
\bottomrule
\end{tabular}
\end{table}

\begin{table}[H]
\centering
\caption{Sesiones, solicitudes, preferencias y distancias.}
\small
\begin{tabular}{@{}p{0.27\linewidth}p{0.69\linewidth}@{}}
\toprule
\textbf{Archivo} & \textbf{Columnas en orden y clave primaria} \\
\midrule
\texttt{horarios.csv} &
\texttt{grupo\_id, dia, inicio, fin}. Clave: los cuatro campos.
El grupo debe existir en \texttt{grupos.csv}. \\
\texttt{solicitudes.csv} &
\texttt{estudiante\_id, asignatura\_id, prioridad}. Clave:
(estudiante, asignatura). Ambos identificadores deben existir. \\
\texttt{preferencias\_}\newline\texttt{profesor.csv} &
\texttt{estudiante\_id, profesor\_id, valor}. Clave:
(estudiante, profesor). Valor entre 0 y 5. \\
\texttt{dias\_libres.csv} &
\texttt{estudiante\_id, dia}. Clave: (estudiante, día). \\
\texttt{companeros.csv} &
\texttt{estudiante\_id, companero\_id}. Clave: (estudiante,
compañero). Ambos son estudiantes distintos. La relación es dirigida. \\
\texttt{distancias.csv} &
\texttt{origen\_id, destino\_id, metros}. Clave: (origen, destino).
Se registra cada par ordenado de salones, incluida la diagonal. \\
\bottomrule
\end{tabular}
\end{table}

Las tablas de preferencias por profesor, días libres y compañeros pueden
contener solo el encabezado. La ausencia de una fila indica que no se
informó esa preferencia. Si un estudiante no define límites de hora, la
exportación registra \texttt{inicio\_min=0}, \texttt{fin\_max=1440}.
Cada estudiante tiene al menos una solicitud y cada grupo al menos una
sesión. Una asignatura puede estar en el catálogo y en una solicitud aunque
carezca de grupos. Esa solicitud podrá quedar pendiente.

\subsection{Ejemplo de filas enlazadas}

Estos fragmentos muestran cómo se reconstruye una parte de la instancia.
El estudiante \texttt{E001} solicita \texttt{A01} con prioridad 5.
El grupo \texttt{A01-G1} pertenece a esa asignatura, admite 12 estudiantes y
ocupa dos sesiones en el mismo salón. El lector reúne las dos franjas bajo
el mismo identificador.

\begin{verbatim}
estudiantes.csv
estudiante_id,inicio_min,fin_max,compacto
E001,480,1080,1

solicitudes.csv
estudiante_id,asignatura_id,prioridad
E001,A01,5

grupos.csv
grupo_id,asignatura_id,profesor_id,salon_id,capacidad
A01-G1,A01,P01,S01,12

horarios.csv
grupo_id,dia,inicio,fin
A01-G1,1,480,570
A01-G1,3,480,570
\end{verbatim}

Los catálogos deben contener \texttt{A01}, \texttt{P01} y \texttt{S01}.
Si la instancia solo tiene el salón \texttt{S01}, \texttt{distancias.csv}
contiene el encabezado \texttt{origen\_id,destino\_id,metros} y la fila
\texttt{S01,S01,0}.

\subsection{Creación, carga y mantenimiento}

\begin{enumerate}
  \item Crear los catálogos con identificadores estables. Para una fuente
        real, convertir horarios a minutos y reemplazar identificadores
        personales por códigos de la instancia antes de guardar.
  \item Crear los grupos con referencias a los catálogos y registrar sus
        franjas en filas separadas. Crear estudiantes, solicitudes y
        preferencias usando esos mismos identificadores.
  \item Ordenar las filas por clave y escribir los once CSV en una carpeta
        nueva. Para datos sintéticos, registrar semilla, escenario, versión
        del esquema y versión de Python en \texttt{meta.json}.
  \item Leer primero los catálogos y estudiantes. Comprobar después las
        referencias al cargar grupos, sesiones, solicitudes, preferencias
        y distancias. Reconstruir una lista de sesiones por grupo y un
        conjunto de solicitudes y preferencias por estudiante.
  \item Validar la instancia reconstruida antes de asignar. Al corregir o
        regenerar una instancia, guardar el conjunto como otra carpeta.
        Así se conservan los datos usados en las mediciones anteriores.
\end{enumerate}

En la implementación de la entrega final, estos comandos generan los archivos
en una carpeta nueva y ejecutan la asignación directamente desde ella.

\begin{verbatim}
python horarios.py generar pequeno ejemplo --formato csv --semilla 202630
python horarios.py resolver ejemplo resultado.json
python horarios.py verificar ejemplo resultado.json
\end{verbatim}

\section{Generación de datos sintéticos}

La instancia se genera con una semilla explícita para que el mismo escenario,
semilla y versión del generador y de Python produzcan los mismos datos. Cada
escenario fija el número de estudiantes, asignaturas, grupos por asignatura,
profesores, salones y solicitudes por estudiante. El generador realiza estos
pasos:

\begin{enumerate}
  \item Crear identificadores de los catálogos. Para los salones se generan
        coordenadas simuladas y se toma la suma de las diferencias absolutas
        por eje. La matriz resultante es simétrica y tiene diagonal cero.
  \item Crear las secciones. Las clases empiezan entre 08:00 y 16:00 en
        pasos de 30 minutos y duran 60, 90 o 120 minutos. Cada sección tiene
        una o dos sesiones en días distintos. Se usan los mismos profesor
        y salón para sus sesiones.
  \item Parametrizar los cupos. Si $R$ es el número total de solicitudes y
        $|G|$ el número de grupos, el cupo base es
        $b=\max(1,\lfloor 0{,}8R/|G|\rfloor)$. El cupo de cada sección
        varía hasta $\lfloor b/3\rfloor$ alrededor de $b$, con mínimo 1.
        Así se puede observar competencia por cupos.
  \item Seleccionar materias sin reemplazo para cada estudiante. Se
        asignan pesos descendentes a los códigos de materia para simular
        demanda desigual. Las prioridades sintéticas van de 1 a 5. Se
        generan valoraciones de profesores, límites de hora, compactación,
        días libres y hasta tres compañeros distintos por estudiante.
\end{enumerate}

La oferta sintética se trata como entrada fija del problema estudiantil.
Para evaluar una situación real se propone recuperar la oferta disponible,
obtener preferencias voluntarias y usar distancias medidas o declaradas con su
procedencia. Los experimentos de esta entrega emplean datos sintéticos y
los casos manuales de referencia.

\section{Escenarios de prueba y salidas esperadas}

\begin{table}[H]
\centering
\caption{Escalas del generador. Se ejecutan las semillas 202630, 202631 y 202632.}
\small
\begin{tabular}{lrrrrrr}
\toprule
Escenario & Estudiantes & Materias & Grupos & Profesores & Salones & Solicitudes \\
\midrule
Pequeño & 5 & 5 & 10 & 3 & 3 & 15 \\
Mediano & 50 & 15 & 45 & 12 & 10 & 250 \\
Grande & 500 & 30 & 120 & 30 & 20 & 3000 \\
\bottomrule
\end{tabular}
\end{table}

El escenario pequeño permite inspeccionar a mano la estructura CSV y revisar
asignaciones concretas. En el mediano se compara la cobertura y la prioridad
atendida sobre las mismas entradas. En el grande se registra además el tiempo
del solucionador. En las tres escalas se comprueba la validez de cada salida
y se conservan las métricas por semilla para reproducir la comparación.

\begin{table}[H]
\centering
\caption{Casos de referencia con resultados comprobables.}
\small
\begin{tabular}{@{}p{0.23\linewidth}p{0.35\linewidth}p{0.38\linewidth}@{}}
\toprule
\textbf{Caso} & \textbf{Configuración} & \textbf{Comprobación esperada} \\
\midrule
Reparación & Una sección de A bloquea ambas opciones de B.
La alternativa de A deja espacio para B. &
Es posible cubrir 2 de 2 solicitudes. Cambiar la sección conserva A. \\
Cupo compartido & Dos estudiantes solicitan la misma asignatura.
Hay una sola sección de cupo 1. &
Máximo de 1 de 2 solicitudes. Entre salidas de igual cobertura,
la referencia atiende prioridad 5. \\
Brecha & Dos estudiantes y una asignatura sin oferta.
Para cubrir tres solicitudes se debe mover al primer estudiante. &
Máximo de 3 de 4 solicitudes. La heurística local puede quedarse en 2. \\
Límite temporal & Dos sesiones del mismo día terminan y empiezan a la
misma hora. Se prueba también un intervalo que invade la otra sesión. &
Se permite la colindancia. Se rechaza la asignación con cruce. \\
CSV inconsistente & Encabezado incorrecto, clave repetida, ID ausente,
franja invertida o distancia sin registrar. &
Se rechaza la entrada antes de construir la asignación. \\
\bottomrule
\end{tabular}
\end{table}

Las salidas esperadas de los tres primeros casos se conservan en
\texttt{salidas\_esperadas.json}. Se contrastan con un oráculo exhaustivo
cuando el número de combinaciones es pequeño. Una referencia se puede construir
sembrando primero una asignación válida y verificando después sus cupos y
franjas. Pueden existir horarios distintos con el mismo valor, por lo que se
compara la validez y los valores de referencia. En instancias mayores se usa
la misma entrada para la línea base y el método propuesto.

\section{Estrategia de validación}

\begin{enumerate}
  \item \textbf{Estructura e integridad de los CSV.} Revisar encabezados,
        número de columnas, valores requeridos, unicidad de claves y
        existencia de referencias. Las filas problemáticas se identifican
        por archivo y línea. Los errores que abarcan varias filas, como una
        matriz de distancias incompleta, se detectan al reconstruir la
        instancia.
  \item \textbf{Dominios y coherencia de la entrada.} Comprobar cupos
        positivos, horarios dentro del dominio, hora de inicio menor que
        la de fin, sesiones de cada grupo sin cruces y matriz de distancias
        completa, simétrica y no negativa, con diagonal cero.
  \item \textbf{Validez de $H$.} Un verificador independiente revisa que
        todos los estudiantes aparezcan en la salida. Después comprueba
        las cuatro postcondiciones mediante la agenda reconstruida,
        las solicitudes y el conteo de cupos. Una asignación parcial puede
        ser válida y las materias sin asignar se reportan como pendientes.
  \item \textbf{Calidad de $H$.} Recalcular cobertura, prioridades, huecos,
        distancias, preferencias de profesor y coincidencias. Comparar
        con casos de referencia o con la línea base según la escala.
        Registrar el tiempo de resolución por separado del tiempo de carga
        del CSV.
\end{enumerate}

\section{Alcance de esta entrega}

Esta primera entrega fija el contrato del problema y describe cómo construir,
intercambiar y verificar las entradas y los casos de referencia. La
implementación del algoritmo y los resultados de ejecución corresponden a
la entrega final.

\section{Conclusión}

El esquema CSV permite construir la misma instancia a partir de un generador
o de fuentes autorizadas y verificar sus relaciones antes de ejecutar el
algoritmo. Los casos de referencia distinguen errores de datos, asignaciones
inválidas y diferencias de calidad entre resultados válidos.

\end{document}
